% !TEX root = ../main.tex
%% Corrigés — Chapitre 4 — Techniques de calcul différentiel et intégral
%% Résultats vérifiés par verifications/chapitre_04.py.

\corrigesChapitre{4}{Techniques de calcul différentiel et intégral}
\corrigesSouples

\begin{blocLivre}{cor-4-1}
\begin{corrige}{1}{Ensembles de définition et transformations}
\demarche{Lister les contraintes : dénominateurs non nuls, arguments de $\sqrt{\ }$ positifs, de $\ln$ strictement positifs.}
\begin{questions}
  \item $\sqrt{x-2}$ : $\res{\intFO{2}{+\infty}}$ ;\quad
    $\dfrac{1}{x^{2}-1}$ : $\res{\R\setminus\ens{-1,1}}$ ;\quad
    $\ln(3-x)$ : $\res{\intOO{-\infty}{3}}$ ;\quad
    $\sqrt{x^{2}-4}$ : $\res{\intOF{-\infty}{-2}\cup\intFO{2}{+\infty}}$.
  \item Translation verticale de $+1$ ; translation horizontale de $-1$
    (la courbe « recule ») ; compression horizontale de facteur $\frac12$ ;
    symétrie par rapport à l'axe des abscisses.
\end{questions}
\end{corrige}
\end{blocLivre}

\begin{blocLivre}{cor-4-2}
\begin{corrige}{2}{Parité, périodicité}
\demarche{Calculer $f(-x)$ et le comparer à $\pm f(x)$ ; pour une somme, prendre un multiple commun des périodes.}
\begin{questions}
  \item $x^{3}-x$ : impaire ; $x^{2}+\cos x$ : paire ; $x+1$ : ni paire ni
    impaire ($f(1)=2$, $f(-1)=0$) ; $x|x|$ : impaire ; $\mathrm{e}^{-x^{2}}$ : paire.
  \item $\cos(3x)$ : période $\dfrac{2\pi}{3}$ ; $\sin\bigl(\frac{x}{2}\bigr)$ :
    période $4\pi$ ; la somme est $\res{4\pi}$-périodique ($4\pi$ est
    multiple commun : $4\pi=6\times\frac{2\pi}{3}$).
  \item Par périodicité, on se ramène à $\intFF{-1}{1}$ ; par parité, à
    $\res{\intFF{0}{1}}$.
\end{questions}
\end{corrige}
\end{blocLivre}

\begin{blocLivre}{cor-4-3}
\begin{corrige}{3}{Calculs de dérivées}
\demarche{Identifier la structure (produit, quotient, composée) avant de dériver.}
\begin{questions}
  \item $3x^{2}-2$ (sur $\R$).
  \item $(x^{2}+2x)\,\mathrm{e}^{x}$ (produit).
  \item $\dfrac{\mathrm{e}^{x}(x+1)-\mathrm{e}^{x}}{(x+1)^{2}}=\dfrac{x\,\mathrm{e}^{x}}{(x+1)^{2}}$
    (sur $\R\setminus\ens{-1}$).
  \item $\dfrac{2x}{1+x^{2}}$ (composée $\ln u$).
  \item $-3\sin x\cos^{2}x$ (composée $u^{3}$).
  \item $\dfrac{x}{\sqrt{x^{2}+1}}$ (composée $\sqrt{u}$).
  \item $2\bigl(1+\tan^{2}(2x)\bigr)$ (sur les intervalles où $\cos(2x)\neq0$).
  \item $\ln x+1-1=\ln x$ (sur $\intOO{0}{+\infty}$) — d'où la primitive
    classique de $\ln$.
\end{questions}
\end{corrige}
\end{blocLivre}

\begin{blocLivre}{cor-4-4}
\begin{corrige}{4}{Variations et extremums}
\demarche{Signe de la dérivée sur chaque intervalle, puis lecture du tableau.}
\begin{questions}
  \item $f'(x)=3x^{2}-3=3(x-1)(x+1)$ :
\begin{center}
\begin{tableau}
  \otTab[2.2]{$x$/0.7, $f'(x)$/0.7, $f$/1.3}{$-\infty$,$-1$,$1$,$+\infty$}
  \tkzTabLine{,+,z,-,z,+,}
  \tkzTabVar{-/$-\infty$, +/$3$, -/$-1$, +/$+\infty$}
\end{tableau}
\end{center}
    Maximum local $f(-1)=\res{3}$, minimum local $f(1)=\res{-1}$.
  \item $g'(x)=(1-x)\,\mathrm{e}^{-x}$ : $g$ croît sur $\intOF{-\infty}{1}$,
    décroît ensuite ; maximum $g(1)=\res{\mathrm{e}^{-1}}$ ; en $+\infty$,
    croissances comparées : $\lim g=0$.
  \item Pour $h(x)=x+\dfrac1x$ : $h'(x)=1-\dfrac{1}{x^{2}}$, négative sur
    $\intOO{0}{1}$, positive ensuite : minimum $h(1)=2$. Donc
    $x+\dfrac1x\geq2$ pour $x>0$, avec égalité uniquement en $\res{x=1}$.
\end{questions}
\end{corrige}
\end{blocLivre}

\begin{blocLivre}{cor-4-5}
\begin{corrige}{5}{Croissances comparées : limites}
\demarche{Reconnaître un rapport de croissances comparées, quitte à poser une variable auxiliaire.}
\begin{questions}
  \item $\res{0}$ (croissance comparée en $0^{+}$).
  \item $\res{0}$ : l'exponentielle l'emporte sur $x^{2}$.
  \item $\res{0}$ : $\ln$ est négligeable devant toute puissance.
  \item $x^{x}=\mathrm{e}^{x\ln x}$ et $x\ln x\to0$ : la limite vaut $\res{1}$.
  \item $\res{+\infty}$.
\end{questions}
\end{corrige}
\end{blocLivre}

\begin{blocLivre}{cor-4-6}
\begin{corrige}{6}{Arcsin, Arccos, Arctan : premiers calculs}
\demarche{Chercher l'angle dans l'intervalle image : $[-\frac{\pi}{2},\frac{\pi}{2}]$, $[0,\pi]$ ou $]-\frac{\pi}{2},\frac{\pi}{2}[$.}
\begin{questions}
  \item $\arcsin\dfrac12=\res{\dfrac{\pi}{6}}$ ;\quad
    $\arccos\Bigl(-\dfrac{\sqrt2}{2}\Bigr)=\res{\dfrac{3\pi}{4}}$ ;\quad
    $\arctan1=\res{\dfrac{\pi}{4}}$ ;\quad $\arctan\sqrt3=\res{\dfrac{\pi}{3}}$.
  \item $\sin\dfrac{5\pi}{6}=\dfrac12$, donc
    $\arcsin\Bigl(\sin\dfrac{5\pi}{6}\Bigr)=\res{\dfrac{\pi}{6}}$
    (et non $\frac{5\pi}{6}$ !) ; de même
    $\arccos\Bigl(\cos\Bigl(-\dfrac{\pi}{3}\Bigr)\Bigr)=\arccos\dfrac12=\res{\dfrac{\pi}{3}}$.
  \item $\bigl(\arctan(2x)\bigr)'=\dfrac{2}{1+4x^{2}}$ sur $\R$ ;
    $\bigl(\arcsin(x^{2})\bigr)'=\dfrac{2x}{\sqrt{1-x^{4}}}$ sur $\intOO{-1}{1}$.
\end{questions}
\end{corrige}
\end{blocLivre}

\begin{blocLivre}{cor-4-7}
\begin{corrige}{7}{Primitives usuelles}
\demarche{Chaque réponse se contrôle en dérivant.}
\begin{questions}
  \item $\dfrac{x^{4}}{4}-\dfrac1x$ (sur $\intOO{0}{+\infty}$ ou $\intOO{-\infty}{0}$).
  \item $\dfrac{\mathrm{e}^{2x}}{2}$ (sur $\R$).
  \item $\dfrac{\sin(3x)}{3}$ (sur $\R$).
  \item $\dfrac12\ln|2x+1|$ (sur un intervalle évitant $-\frac12$).
  \item $\dfrac12\ln(x^{2}+1)$ — forme $\dfrac{u'}{u}$ (sur $\R$).
  \item $\dfrac{\sin^{2}x}{2}$ — forme $u'u$ (sur $\R$).
  \item $\dfrac12\arctan(2x)$ (sur $\R$).
  \item $\arcsin x$ (sur $\intOO{-1}{1}$).
\end{questions}
\end{corrige}
\end{blocLivre}

\begin{blocLivre}{cor-4-8}
\begin{corrige}{8}{Équations homogènes du premier ordre}
\demarche{Les solutions de $y'+ay=0$ sont les $C\,\mathrm{e}^{-ax}$ ; la condition initiale fixe $C$.}
\begin{questions}
  \item Solutions : $y=C\,\mathrm{e}^{-3x}$, $C\in\R$ ; la condition $y(0)=2$
    donne $\res{y=2\,\mathrm{e}^{-3x}}$.
  \item Solutions de $y'-2y=0$ : $C\,\mathrm{e}^{2x}$ ; d'où $\res{y=-\mathrm{e}^{2x}}$.
  \item Si $y(x_0)=0$, alors $C\,\mathrm{e}^{-ax_0}=0$ donc $C=0$ : la solution
    est la fonction nulle (c'est l'unicité du problème de Cauchy).
\end{questions}
\end{corrige}
\end{blocLivre}

\begin{blocLivre}{cor-4-9}
\begin{corrige}{9}{Deux études de fonctions}
\demarche{Tableau de variations, puis lecture des bornes atteintes.}
\begin{questions}
  \item $f$ est impaire, et $f'(x)=\dfrac{1-x^{2}}{(1+x^{2})^{2}}$ :
\begin{center}
\begin{tableau}
  \otTab[2.2]{$x$/0.7, $f'(x)$/0.7, $f$/1.3}{$-\infty$,$-1$,$1$,$+\infty$}
  \tkzTabLine{,-,z,+,z,-,}
  \tkzTabVar{+/$0$, -/$-\frac12$, +/$\frac12$, -/$0$}
\end{tableau}
\end{center}
    Donc $f$ est bornée : $\res{\max f=f(1)=\tfrac12}$ et
    $\res{\min f=f(-1)=-\tfrac12}$.
  \item $g'(x)=\dfrac{1-\ln x}{x^{2}}$ : $g$ croît sur $\intOF{0}{\mathrm{e}}$,
    décroît ensuite ; maximum $g(\mathrm{e})=\dfrac{1}{\mathrm{e}}$ ;
    $\lim_{0^{+}}g=-\infty$, $\lim_{+\infty}g=0$. Comme le maximum est
    atteint \emph{uniquement} en $\mathrm{e}$ : $g(\mathrm{e})>g(\pi)$, soit
    $\dfrac{1}{\mathrm{e}}>\dfrac{\ln\pi}{\pi}$, d'où $\pi>\mathrm{e}\ln\pi=\ln(\pi^{\mathrm{e}})$,
    et en composant par l'exponentielle (croissante) :
    $\res{\mathrm{e}^{\pi}>\pi^{\mathrm{e}}}$.
\end{questions}
\end{corrige}
\end{blocLivre}

\begin{blocLivre}{cor-4-10}
\begin{corrige}{10}{Des inégalités par étude de fonctions}
\demarche{Étudier la différence des deux membres : son minimum donne l'inégalité.}
\begin{questions}
  \item Pour $h(x)=\mathrm{e}^{x}-1-x$ : $h'(x)=\mathrm{e}^{x}-1$, du signe de $x$ ;
    $h$ décroît puis croît, avec $h(0)=0$ : $h\geq0$ partout, nulle
    seulement en $0$.
  \item Pour $x>-1$, posons $u=\ln(1+x)$ : la question 1 donne
    $\mathrm{e}^{u}\geq1+u$, c'est-à-dire $1+x\geq1+\ln(1+x)$, d'où
    $\ln(1+x)\leq x$. (Ou : étude directe de $x\mapsto x-\ln(1+x)$.)
  \item Pour $k(x)=\ln(1+x)-x+\dfrac{x^{2}}{2}$ :
    $k'(x)=\dfrac{1}{1+x}-1+x=\dfrac{x^{2}}{1+x}\geq0$ sur $\intFO{0}{+\infty}$,
    et $k(0)=0$ : donc $k\geq0$, ce qui est l'inégalité voulue.
\end{questions}
\end{corrige}
\end{blocLivre}

\begin{blocLivre}{cor-4-11}
\begin{corrige}{11}{Dérivée d'une fonction réciproque}
\demarche{La pente de la réciproque est l'inverse de la pente, lue au point symétrique.}
\begin{questions}
  \item $f'(x)=1+\mathrm{e}^{x}>0$ : $f$ est strictement croissante et continue,
    avec $\lim_{-\infty}f=-\infty$ et $\lim_{+\infty}f=+\infty$ : $f$ est
    une bijection de $\R$ sur $\R$.
  \item $f'$ ne s'annule pas, donc $f^{-1}$ est dérivable. Comme $f(0)=1$ :
    \[
      \bigl(f^{-1}\bigr)'(1)=\frac{1}{f'(0)}=\frac{1}{1+\mathrm{e}^{0}}=\res{\frac12}.
    \]
  \item $f^{-1}(1)=0$, donc la tangente est
    $\res{y=\dfrac{x-1}{2}}$.
\end{questions}
\end{corrige}
\end{blocLivre}

\begin{blocLivre}{cor-4-12}
\begin{corrige}{12}{Identités avec Arcsin, Arccos, Arctan}
\demarche{Soit on compose des formules de trigonométrie, soit on dérive : une fonction de dérivée nulle sur un intervalle est constante.}
\begin{questions}
  \item Posons $t=\arcsin x\in\bigl[-\frac{\pi}{2},\frac{\pi}{2}\bigr]$ :
    $\cos t\geq0$, donc $\cos t=\sqrt{1-\sin^{2}t}=\sqrt{1-x^{2}}$.
  \item La fonction $x\mapsto\arcsin x+\arccos x$ est continue sur
    $\intFF{-1}{1}$, dérivable sur $\intOO{-1}{1}$ de dérivée
    $\dfrac{1}{\sqrt{1-x^{2}}}-\dfrac{1}{\sqrt{1-x^{2}}}=0$ : elle est
    constante, égale à sa valeur en $0$ :
    $\arcsin0+\arccos0=\dfrac{\pi}{2}$.
  \item Sur $\intOO{0}{+\infty}$, la dérivée de
    $x\mapsto\arctan x+\arctan\dfrac1x$ vaut
    $\dfrac{1}{1+x^{2}}+\dfrac{-1/x^{2}}{1+1/x^{2}}=\dfrac{1}{1+x^{2}}-\dfrac{1}{x^{2}+1}=0$ :
    la fonction est constante, égale à sa valeur en $1$, soit
    $2\arctan1=\dfrac{\pi}{2}$. Par imparité, elle vaut $-\dfrac{\pi}{2}$
    pour $x<0$.
  \item Posons $t=\arctan x$ : $\sin(2t)=2\sin t\cos t=2\tan t\cos^{2}t
    =\dfrac{2\tan t}{1+\tan^{2}t}=\dfrac{2x}{1+x^{2}}$.
\end{questions}
\end{corrige}
\end{blocLivre}

\begin{blocLivre}{cor-4-13}
\begin{corrige}{13}{Fonctions hyperboliques}
\demarche{Revenir aux exponentielles : toutes les identités s'y ramènent.}
\begin{questions}
  \item $\operatorname{ch}^{2}x-\operatorname{sh}^{2}x
    =\dfrac{(\mathrm{e}^{x}+\mathrm{e}^{-x})^{2}-(\mathrm{e}^{x}-\mathrm{e}^{-x})^{2}}{4}
    =\dfrac{4\,\mathrm{e}^{x}\mathrm{e}^{-x}}{4}=1$.
  \item En développant
    $\operatorname{sh}a\operatorname{ch}b+\operatorname{ch}a\operatorname{sh}b
    =\dfrac{(\mathrm{e}^{a}-\mathrm{e}^{-a})(\mathrm{e}^{b}+\mathrm{e}^{-b})
    +(\mathrm{e}^{a}+\mathrm{e}^{-a})(\mathrm{e}^{b}-\mathrm{e}^{-b})}{4}
    =\dfrac{2\,\mathrm{e}^{a+b}-2\,\mathrm{e}^{-(a+b)}}{4}=\operatorname{sh}(a+b)$.
  \item $\operatorname{th}'=\dfrac{\operatorname{ch}^{2}-\operatorname{sh}^{2}}{\operatorname{ch}^{2}}
    =1-\operatorname{th}^{2}=\dfrac{1}{\operatorname{ch}^{2}}>0$ :
    $\operatorname{th}$ est strictement croissante, impaire, avec
    $\lim_{\pm\infty}\operatorname{th}=\pm1$.
  \item $\operatorname{ch}x=2$ équivaut à $\mathrm{e}^{x}+\mathrm{e}^{-x}=4$, soit, en
    posant $X=\mathrm{e}^{x}>0$ : $X^{2}-4X+1=0$, d'où $X=2\pm\sqrt3$. Les deux
    valeurs sont positives et inverses l'une de l'autre :
    $\res{x=\pm\ln\bigl(2+\sqrt3\bigr)}$.
\end{questions}
\end{corrige}
\end{blocLivre}

\begin{blocLivre}{cor-4-14}
\begin{corrige}{14}{Primitives : formes composées}
\demarche{Repérer $u$ et $u'$ ; vérifier ensuite en dérivant.}
\begin{questions}
  \item $\dfrac{\mathrm{e}^{x^{2}}}{2}$ — forme $u'\mathrm{e}^{u}$ (sur $\R$).
  \item $\tan x=\dfrac{\sin x}{\cos x}=-\dfrac{u'}{u}$ avec $u=\cos$ :
    primitive $-\ln|\cos x|$ (sur un intervalle où $\cos\neq0$).
  \item $\dfrac{1/x}{\ln x}$ : primitive $\ln|\ln x|$ (sur $\intOO{1}{+\infty}$
    par exemple).
  \item $\dfrac{\sin^{4}x}{4}$ — forme $u'u^{3}$.
  \item $\sqrt{1+x^{2}}$ — forme $\dfrac{u'}{2\sqrt{u}}$ (sur $\R$).
  \item $\dfrac{\mathrm{e}^{x}}{1+(\mathrm{e}^{x})^{2}}$ — forme $\dfrac{u'}{1+u^{2}}$ :
    primitive $\arctan(\mathrm{e}^{x})$ (sur $\R$).
\end{questions}
\end{corrige}
\end{blocLivre}

\begin{blocLivre}{cor-4-15}
\begin{corrige}{15}{Intégration par parties}
\demarche{Dériver le polynôme, primitiver l'exponentielle ou le sinus ; pour $\ln$, dériver $\ln$.}
\begin{questions}
  \item $\displaystyle\int_{0}^{1}x\,\mathrm{e}^{x}\mathrm{d}x
    =\bigl[x\,\mathrm{e}^{x}\bigr]_{0}^{1}-\int_{0}^{1}\mathrm{e}^{x}\mathrm{d}x
    =\mathrm{e}-(\mathrm{e}-1)=\res{1}$.
  \item $\displaystyle\int_{1}^{\mathrm{e}}\ln t\,\mathrm{d}t
    =\bigl[t\ln t\bigr]_{1}^{\mathrm{e}}-\int_{1}^{\mathrm{e}}1\,\mathrm{d}t
    =\mathrm{e}-(\mathrm{e}-1)=\res{1}$.
  \item $\displaystyle\int_{0}^{\pi/2}t\sin t\,\mathrm{d}t
    =\bigl[-t\cos t\bigr]_{0}^{\pi/2}+\int_{0}^{\pi/2}\cos t\,\mathrm{d}t
    =0+1=\res{1}$.
  \item Deux intégrations par parties successives :
    $\displaystyle\int_{0}^{1}x^{2}\mathrm{e}^{x}\mathrm{d}x
    =\bigl[x^{2}\mathrm{e}^{x}\bigr]_{0}^{1}-2\int_{0}^{1}x\,\mathrm{e}^{x}\mathrm{d}x
    =\mathrm{e}-2=\res{\mathrm{e}-2}$ (avec la question 1).
\end{questions}
\end{corrige}
\end{blocLivre}

\begin{blocLivre}{cor-4-16}
\begin{corrige}{16}{Changements de variable}
\demarche{Changer la variable, la différentielle et les bornes — les trois en même temps.}
\begin{questions}
  \item Avec $u=\sqrt{x}$, $x=u^{2}$, $\mathrm{d}x=2u\,\mathrm{d}u$ :
    \[
      \int_{0}^{1}\frac{\mathrm{d}x}{1+\sqrt{x}}
      =\int_{0}^{1}\frac{2u}{1+u}\,\mathrm{d}u
      =2\int_{0}^{1}\Bigl(1-\frac{1}{1+u}\Bigr)\mathrm{d}u
      =\res{2-2\ln2}.
    \]
  \item Avec $u=\mathrm{e}^{t}$, $\mathrm{d}u=u\,\mathrm{d}t$ :
    \[
      \int_{0}^{1}\frac{\mathrm{d}t}{\mathrm{e}^{t}+1}
      =\int_{1}^{\mathrm{e}}\frac{\mathrm{d}u}{u(u+1)}
      =\int_{1}^{\mathrm{e}}\Bigl(\frac1u-\frac{1}{u+1}\Bigr)\mathrm{d}u
      =\Bigl[\ln\frac{u}{u+1}\Bigr]_{1}^{\mathrm{e}}
      =\res{1+\ln2-\ln(\mathrm{e}+1)}.
    \]
  \item Avec $u=\sin t$, $\mathrm{d}u=\cos t\,\mathrm{d}t$ et
    $\cos^{2}t=1-u^{2}$ :
    \[
      \int_{0}^{\pi/2}\cos^{3}t\,\mathrm{d}t=\int_{0}^{1}(1-u^{2})\,\mathrm{d}u
      =1-\frac13=\res{\frac23}.
    \]
\end{questions}
\end{corrige}
\end{blocLivre}

\begin{blocLivre}{cor-4-17}
\begin{corrige}{17}{Premier ordre avec second membre}
\demarche{Solution particulière + solutions de l'homogène ; la condition initiale vient en dernier.}
\begin{questions}
  \item Homogène : $C\,\mathrm{e}^{-2x}$. On cherche $y_p=\mu\,\mathrm{e}^{3x}$ :
    $3\mu+2\mu=1$ donne $\mu=\dfrac15$. Solutions :
    $y=\dfrac{\mathrm{e}^{3x}}{5}+C\,\mathrm{e}^{-2x}$ ; la condition $y(0)=0$ donne
    $C=-\dfrac15$ :
    $\res{y=\dfrac{\mathrm{e}^{3x}-\mathrm{e}^{-2x}}{5}}$.
  \item $y_p'=-\alpha\sin x+\beta\cos x$, et
    $y_p'-y_p=(\beta-\alpha)\cos x-(\alpha+\beta)\sin x=\cos x$ impose
    $\beta-\alpha=1$ et $\alpha+\beta=0$ : $\alpha=-\dfrac12$,
    $\beta=\dfrac12$, d'où $y_p=\dfrac{\sin x-\cos x}{2}$ et les solutions
    $\res{y=\dfrac{\sin x-\cos x}{2}+C\,\mathrm{e}^{x}}$.
  \item Pour le second membre $1$ : $y_{p_1}=1$. Pour $\mathrm{e}^{-x}$, qui
    résout l'homogène, on cherche $y_{p_2}=\mu x\,\mathrm{e}^{-x}$ :
    $y_{p_2}'+y_{p_2}=\mu\,\mathrm{e}^{-x}$, donc $\mu=1$. Par superposition :
    $\res{y=1+x\,\mathrm{e}^{-x}+C\,\mathrm{e}^{-x}}$.
\end{questions}
\end{corrige}
\end{blocLivre}

\begin{blocLivre}{cor-4-18}
\begin{corrige}{18}{Variation de la constante}
\demarche{Faire varier la constante de l'homogène : il ne reste qu'une primitive à calculer.}
\begin{questions}
  \item Homogène : $y'=-\dfrac{y}{x}$, de solutions $\dfrac{C}{x}$. On pose
    $y=\dfrac{C(x)}{x}$ : l'équation donne $\dfrac{C'(x)}{x}=3x$, soit
    $C'(x)=3x^{2}$, d'où $C(x)=x^{3}$ et $y_p=x^{2}$. Solutions sur
    $\intOO{0}{+\infty}$ : $y=x^{2}+\dfrac{C}{x}$ ; la condition $y(1)=2$
    donne $C=1$ : $\res{y=x^{2}+\dfrac1x}$.
  \item Homogène : $C\,\mathrm{e}^{x^{2}}$ (car $-\!\int(-2x)\,\mathrm{d}x=x^{2}$). En posant
    $y=C(x)\,\mathrm{e}^{x^{2}}$ : $C'(x)\,\mathrm{e}^{x^{2}}=\mathrm{e}^{x^{2}}$, donc
    $C'(x)=1$ et $C(x)=x$. Solutions :
    $\res{y=(x+C)\,\mathrm{e}^{x^{2}},\ C\in\R}$.
\end{questions}
\end{corrige}
\end{blocLivre}

\begin{blocLivre}{cor-4-19}
\begin{corrige}{19}{Second ordre à coefficients constants}
\demarche{Équation caractéristique, forme des solutions selon le discriminant, second membre à la fin.}
\begin{questions}
  \item $y''+\omega^{2}y=0$ : racines $\pm\mathrm{i}\omega$, solutions
    $A\cos(\omega x)+B\sin(\omega x)$.\\
    $y''-3y'+2y=0$ : racines $1$ et $2$, solutions $A\,\mathrm{e}^{x}+B\,\mathrm{e}^{2x}$.\\
    $y''+4y'+4y=0$ : racine double $-2$, solutions $(A+Bx)\,\mathrm{e}^{-2x}$.\\
    $y''+2y'+5y=0$ : racines $-1\pm2\mathrm{i}$, solutions
    $\mathrm{e}^{-x}\bigl(A\cos(2x)+B\sin(2x)\bigr)$.
  \item On cherche $y_p=\mu\cos(2x)$ (par parité du second membre) :
    $-4\mu+\mu=1$ donne $\mu=-\dfrac13$. Solutions :
    $\res{y=-\dfrac{\cos(2x)}{3}+A\cos x+B\sin x}$.
  \item $\lambda=1$ est racine simple de $r^{2}-3r+2$ : on pose
    $y_p=\mu x\,\mathrm{e}^{x}$. Alors
    $y_p''-3y_p'+2y_p=\mu\,\mathrm{e}^{x}\bigl((2+x)-3(1+x)+2x\bigr)=-\mu\,\mathrm{e}^{x}$,
    d'où $\mu=-1$ : $\res{y=-x\,\mathrm{e}^{x}+A\,\mathrm{e}^{x}+B\,\mathrm{e}^{2x}}$.
  \item $y=A\cos(\omega x)+B\sin(\omega x)$ avec $y(0)=A=y_0$ et
    $y'(0)=B\omega=v_0$ :
    $y=y_0\cos(\omega x)+\dfrac{v_0}{\omega}\sin(\omega x)$, c'est-à-dire
    $A\cos(\omega x-\varphi)$ avec $A=\sqrt{y_0^{2}+(v_0/\omega)^{2}}$ —
    l'oscillateur harmonique, d'amplitude $A$ et de phase $\varphi$.
\end{questions}
\end{corrige}
\end{blocLivre}

\begin{blocLivre}{cor-4-20}
\begin{corrige}{20}{La formule de Machin}
\demarche{Comparer des tangentes ne suffit pas : il faut aussi encadrer les angles.}
\begin{questions}
  \item Posons $S=\arctan\dfrac12+\arctan\dfrac13$. Chaque terme est dans
    $\bigl]0,\frac{\pi}{4}\bigr]$, donc $S\in\bigl]0,\frac{\pi}{2}\bigr[$, et
    \[
      \tan S=\frac{\frac12+\frac13}{1-\frac12\cdot\frac13}=\frac{5/6}{5/6}=1 :
    \]
    le seul angle de $\bigl]0,\frac{\pi}{2}\bigr[$ de tangente $1$ est
    $\res{\dfrac{\pi}{4}}$.
  \item $\tan(2a)=\dfrac{2/5}{1-1/25}=\res{\dfrac{5}{12}}$, puis
    $\tan(4a)=\dfrac{2\cdot5/12}{1-25/144}=\res{\dfrac{120}{119}}$.
  \item Comme $\tan a=\dfrac15<\tan\dfrac{\pi}{8}$, on a $a<\dfrac{\pi}{8}$
    et $4a\in\bigl]0,\frac{\pi}{2}\bigr[$, donc
    $4a-\dfrac{\pi}{4}\in\bigl]-\frac{\pi}{4},\frac{\pi}{4}\bigr[$. Sa tangente vaut
    \[
      \frac{\frac{120}{119}-1}{1+\frac{120}{119}}=\frac{1/119}{239/119}=\frac{1}{239},
    \]
    donc $4a-\dfrac{\pi}{4}=\arctan\dfrac{1}{239}$, ce qui est la formule de
    Machin.
\end{questions}
\end{corrige}
\end{blocLivre}

\begin{blocLivre}{cor-4-21}
\begin{corrige}{21}{Un encadrement de Arctan}
\demarche{Encore des différences à étudier : leurs dérivées sont des fractions rationnelles simples.}
\begin{questions}
  \item Pour $u(x)=x-\arctan x$ : $u'(x)=1-\dfrac{1}{1+x^{2}}
    =\dfrac{x^{2}}{1+x^{2}}\geq0$ et $u(0)=0$ : $u\geq0$ sur $\intFO{0}{+\infty}$.
  \item Pour $v(x)=\arctan x-x+\dfrac{x^{3}}{3}$ :
    $v'(x)=\dfrac{1}{1+x^{2}}-1+x^{2}=\dfrac{x^{4}}{1+x^{2}}\geq0$ et
    $v(0)=0$ : $v\geq0$ sur $\intFO{0}{+\infty}$.
  \item Pour $x>0$ : $1-\dfrac{x^{2}}{3}\leq\dfrac{\arctan x}{x}\leq1$, et le
    théorème d'encadrement donne
    $\res{\displaystyle\lim_{x\to0}\frac{\arctan x}{x}=1}$ (les deux membres
    étant pairs, la limite vaut aussi à gauche).
\end{questions}
\end{corrige}
\end{blocLivre}

\begin{blocLivre}{cor-4-22}
\begin{corrige}{22}{Les intégrales de Wallis}
\demarche{Une intégration par parties bien choisie produit une relation de récurrence.}
\begin{questions}
  \item $I_0=\dfrac{\pi}{2}$ et $I_1=\bigl[-\cos t\bigr]_0^{\pi/2}=1$.
  \item Sur $\bigl[0,\frac{\pi}{2}\bigr]$ : $0\leq\sin t\leq1$, donc
    $\sin^{n+1}t\leq\sin^{n}t$, d'où $I_{n+1}\leq I_n$ ; et $\sin^{n}t>0$
    sur $\bigl]0,\frac{\pi}{2}\bigr[$ donne $I_n>0$.
  \item On intègre par parties $I_{n+2}=\displaystyle\int_0^{\pi/2}\sin t\cdot\sin^{n+1}t\,\mathrm{d}t$,
    en primitivant $\sin t$ :
    \[
      I_{n+2}=\bigl[-\cos t\sin^{n+1}t\bigr]_0^{\pi/2}
      +(n+1)\int_0^{\pi/2}\cos^{2}t\,\sin^{n}t\,\mathrm{d}t
      =(n+1)\,(I_n-I_{n+2}),
    \]
    car $\cos^{2}t=1-\sin^{2}t$. D'où $(n+2)\,I_{n+2}=(n+1)\,I_n$.
  \item $I_2=\dfrac12\cdot\dfrac{\pi}{2}=\res{\dfrac{\pi}{4}}$ ;\quad
    $I_3=\dfrac23\,I_1=\res{\dfrac23}$ ;\quad
    $I_4=\dfrac34\,I_2=\res{\dfrac{3\pi}{16}}$.
\end{questions}
\end{corrige}
\end{blocLivre}

\begin{blocLivre}{cor-4-23}
\begin{corrige}{23}{Refroidissement de Newton}
\demarche{Une équation du premier ordre à coefficients constants, pilotée par l'écart à la température ambiante.}
\begin{questions}
  \item L'équation s'écrit $T'+kT=20k$ ; la constante $T=20$ est solution
    particulière, donc $T(t)=20+C\,\mathrm{e}^{-kt}$ ; avec $T(0)=90$ :
    $\res{T(t)=20+70\,\mathrm{e}^{-kt}}$.
  \item $T(10)=60$ donne $\mathrm{e}^{-10k}=\dfrac{40}{70}=\dfrac47$, d'où
    $\res{k=\dfrac{1}{10}\ln\dfrac74}$.
  \item $T(t)=30$ équivaut à $\mathrm{e}^{-kt}=\dfrac{10}{70}=\dfrac17$, d'où
    $t=\dfrac{\ln7}{k}=\res{\dfrac{10\ln7}{\ln(7/4)}\approx35\ \text{min}}$
    (précisément $34{,}8$ min).
  \item $T(t)-20=70\,\mathrm{e}^{-kt}>0$ pour tout $t$ : la température reste
    strictement supérieure à l'ambiante, dont elle s'approche
    asymptotiquement.
\end{questions}
\end{corrige}
\end{blocLivre}

\begin{blocLivre}{cor-4-24}
\begin{corrige}{24}{Équations fonctionnelles et équations différentielles}
\demarche{Transformer l'équation fonctionnelle en équation différentielle, puis invoquer l'unicité de Cauchy.}
\begin{questions}
  \item Avec $x=y=0$ : $f(0)=f(0)^{2}$, donc $f(0)\in\ens{0,1}$. Si
    $f(0)=0$, alors $f(x)=f(x)f(0)=0$ pour tout $x$ : exclu. Donc $f(0)=1$.
    Si $f(a)=0$ pour un $a$, alors $f(x)=f(x-a)f(a)=0$ pour tout $x$ :
    exclu, donc $f$ ne s'annule pas. Enfin
    $f(x)=f\bigl(\frac{x}{2}\bigr)^{2}>0$.
  \item À $x$ fixé, dérivons $y\mapsto f(x+y)=f(x)f(y)$ :
    $f'(x+y)=f(x)f'(y)$ ; en $y=0$ : $f'(x)=f'(0)\,f(x)$.
  \item En posant $c=f'(0)$ : $f$ résout $y'=cy$ avec $y(0)=1$, donc
    $f(x)=\mathrm{e}^{cx}$. Réciproquement, toute fonction
    $x\mapsto\mathrm{e}^{cx}$ convient : les solutions sont exactement les
    $\res{f:x\mapsto\mathrm{e}^{cx},\ c\in\R}$.
  \item Équation caractéristique $r^{2}-1=0$ : solutions
    $y=A\,\mathrm{e}^{x}+B\,\mathrm{e}^{-x}$.
  \item $y(0)=A+B=1$ et $y'(0)=A-B=0$ donnent $A=B=\dfrac12$ :
    l'unique solution est
    $\res{y=\dfrac{\mathrm{e}^{x}+\mathrm{e}^{-x}}{2}=\operatorname{ch}x}$.
  \item De même, $A+B=0$ et $A-B=1$ donnent $A=-B=\dfrac12$ :
    $\res{y=\operatorname{sh}x}$. Les fonctions $\operatorname{ch}$ et
    $\operatorname{sh}$ sont ainsi \emph{caractérisées} par l'équation
    $y''=y$ et leurs conditions initiales.
\end{questions}
\end{corrige}
\end{blocLivre}

% !TEX root = ../../main.tex
%% Corrigé du défi du chapitre 4 — vérifié par verifications/defis_01_05.py.

\begin{blocLivre}{cor-4-25a}
\begin{corrige}{25}{Oscillateur amorti, régime forcé et résonance}
\demarche{Équation caractéristique pour le régime libre, exponentielle
complexe pour le régime forcé, étude d'un trinôme pour la résonance.}
\partieExo{A}{Le régime libre ($f=0$)}
\begin{questions}
  \item Équation caractéristique $r^{2}+2\lambda r+\omega_0^{2}=0$, de
    discriminant réduit $\lambda^{2}-\omega_0^{2}$.
    \begin{itemize}[label={\color{otNavy}\textbullet},leftmargin=1.2em,itemsep=0.4mm,topsep=0.6mm]
      \item $\lambda>\omega_0$ (apériodique) : deux racines réelles
        $r_\pm=-\lambda\pm\sqrt{\lambda^{2}-\omega_0^{2}}$,
        $y=C_1\mathrm{e}^{r_+t}+C_2\mathrm{e}^{r_-t}$ ;
      \item $\lambda=\omega_0$ (critique) : racine double $-\lambda$,
        $y=(C_1+C_2t)\,\mathrm{e}^{-\lambda t}$ ;
      \item $\lambda<\omega_0$ (pseudo-périodique) : racines
        $-\lambda\pm\mathrm{i}\Omega$,
        $y=\mathrm{e}^{-\lambda t}\bigl(C_1\cos\Omega t+C_2\sin\Omega t\bigr)$.
    \end{itemize}
  \item $(2,\sqrt3)$ : racines $-1$, $-3$ ; $C_1+C_2=1$ et
    $-C_1-3C_2=0$ donnent
    $\res{y=\tfrac32\mathrm{e}^{-t}-\tfrac12\mathrm{e}^{-3t}}$.
    $(1,1)$ : $\res{y=(1+t)\,\mathrm{e}^{-t}}$.
    $(1,2)$ : $\Omega=\sqrt3$, $C_1=1$, $-C_1+\sqrt3\,C_2=0$, d'où
    $\res{y=\mathrm{e}^{-t}\bigl(\cos\sqrt3\,t+\tfrac{1}{\sqrt3}\sin\sqrt3\,t\bigr)}$.
  \item Dans le cas apériodique, $r_-<r_+<0$ car
    $\sqrt{\lambda^{2}-\omega_0^{2}}<\lambda$ ; dans les deux autres,
    $t^{k}\mathrm{e}^{-\lambda t}\to0$ (croissances comparées) et
    $\cos$, $\sin$ sont bornées. Toute solution tend vers
    \res{0} : le régime libre s'éteint.
\end{questions}
\begin{center}
\begin{tikzpicture}[x=1.2cm,y=2.4cm,line cap=round]
  \draw[otRule,line width=0.3pt] (0,0.5) -- (6.2,0.5) (0,1) -- (6.2,1);
  \draw[otMuted,line width=0.5pt,-{Stealth[length=1.6mm]}] (0,0) -- (6.4,0);
  \draw[otMuted,line width=0.5pt,-{Stealth[length=1.6mm]}] (0,-0.2) -- (0,1.2);
  \draw[otNavy!55,line width=1pt] plot[domain=0:6.2,samples=120]
    ({\x},{1.5*exp(-\x)-0.5*exp(-3*\x)});
  \draw[otGreenDark,line width=1.1pt] plot[domain=0:6.2,samples=120]
    ({\x},{(1+\x)*exp(-\x)});
  \draw[otNavy,line width=1.2pt] plot[domain=0:6.2,samples=200]
    ({\x},{exp(-\x)*(cos(deg(sqrt(3)*\x))+sin(deg(sqrt(3)*\x))/sqrt(3))});
  \node[font=\scriptsize,text=otNavy!70,anchor=west] at (2.3,0.42) {apériodique};
  \node[font=\scriptsize,text=otGreenDark,anchor=west] at (1.6,0.68) {critique};
  \node[font=\scriptsize,text=otNavy,anchor=north west] at (1.9,-0.07) {pseudo-périodique};
  \node[font=\tiny,text=otMuted,anchor=east] at (0,1) {$1$};
  \node[font=\scriptsize,text=otMuted,below] at (6.3,0) {$t$};
\end{tikzpicture}\\[0.5mm]
\legende{Les trois régimes libres ($y(0)=1$, $y'(0)=0$) : le régime
critique est celui qui revient le plus vite à l'équilibre sans osciller.}
\end{center}
\end{corrige}
\end{blocLivre}

\begin{blocLivre}{cor-4-25b}
\begin{corrigeSuite}
\partieExo{B}{Le régime forcé ($f(t)=\cos\omega t$)}
\begin{questions}[start=4]
  \item Pour $z=Z\mathrm{e}^{\mathrm{i}\omega t}$ :
    $z''+2\lambda z'+\omega_0^{2}z
    =\bigl(-\omega^{2}+2\mathrm{i}\lambda\omega+\omega_0^{2}\bigr)Z\mathrm{e}^{\mathrm{i}\omega t}$,
    et le facteur est non nul (partie imaginaire $2\lambda\omega\neq0$) :
    $\res{Z=\dfrac{1}{\omega_0^{2}-\omega^{2}+2\mathrm{i}\lambda\omega}}$
    convient.
  \item Les coefficients étant réels, $y_p=\operatorname{Re}z$ est
    solution de $(E_f)$ avec $f=\operatorname{Re}\mathrm{e}^{\mathrm{i}\omega t}=\cos\omega t$.
    Écrivons $Z=A\,\mathrm{e}^{-\mathrm{i}\varphi}$ avec $A=|Z|=A(\omega)$ :
    $\mathrm{e}^{\mathrm{i}\varphi}=A\bigl(\omega_0^{2}-\omega^{2}+2\mathrm{i}\lambda\omega\bigr)$,
    soit
    \[
    \cos\varphi=(\omega_0^{2}-\omega^{2})\,A(\omega),\qquad
    \sin\varphi=2\lambda\omega\,A(\omega)>0 ,
    \]
    d'où $\varphi\in\intOO{0}{\pi}$, et
    $y_p(t)=\operatorname{Re}\bigl(A\mathrm{e}^{\mathrm{i}(\omega t-\varphi)}\bigr)
    =\res{A(\omega)\cos(\omega t-\varphi)}$.
  \item Par linéarité, $y-y_p$ est solution de $(E_0)$, donc tend vers
    $0$ (question 3) : $\res{y(t)-y_p(t)\to0}$. La réponse « oublie » les
    conditions initiales et ne garde que la pulsation imposée.
\end{questions}
\partieExo{C}{La résonance}
\begin{questions}[start=7]
  \item Avec $u=\omega^{2}$ :
    $D(u)=(\omega_0^{2}-u)^{2}+4\lambda^{2}u=u^{2}-2(\omega_0^{2}-2\lambda^{2})u+\omega_0^{4}$,
    trinôme minimal en $u^{*}=\omega_0^{2}-2\lambda^{2}$. Si
    $\lambda<\frac{\omega_0}{\sqrt2}$, $u^{*}>0$ : $D$ décroît puis croît
    sur $\R_{+}$, et $A=D^{-1/2}$ est maximale en
    $\res{\omega_r=\sqrt{\omega_0^{2}-2\lambda^{2}}}$, avec
    $D(u^{*})=\omega_0^{4}-(\omega_0^{2}-2\lambda^{2})^{2}
    =4\lambda^{2}(\omega_0^{2}-\lambda^{2})$, d'où
    $\res{A_{\max}=\dfrac{1}{2\lambda\sqrt{\omega_0^{2}-\lambda^{2}}}}$.
    Si $\lambda\geqslant\frac{\omega_0}{\sqrt2}$, $u^{*}\leqslant0$ : $D$
    croît sur $\R_{+}$, $A$ décroît, pas de résonance.
  \item $A(0)=\dfrac{1}{\omega_0^{2}}$ et $A(\omega_0)=\dfrac{1}{2\lambda\omega_0}$,
    d'où $\res{\dfrac{A(\omega_0)}{A(0)}=\dfrac{\omega_0}{2\lambda}}$.
    Quand $\lambda\to0$ : $\omega_r\to\omega_0$ et
    $A_{\max}\sim\dfrac{1}{2\lambda\omega_0}\to+\infty$. Un système peu
    amorti, excité près de sa pulsation propre, répond avec une amplitude
    considérable : c'est pourquoi on amortit les structures (ponts,
    machines, suspensions) soumises à des sollicitations périodiques. La
    figure de l'énoncé montre le pic s'affiner et monter quand $\lambda$
    diminue, et disparaître pour $\lambda=0{,}8>\frac{1}{\sqrt2}$.
\end{questions}
\end{corrigeSuite}
\end{blocLivre}

